% Part: first-order-logic
% Chapter: tableaux
% Section: derivations

\documentclass[../../../include/open-logic-section]{subfiles}

\begin{document}

\iftag{FOL}
      {\olfileid{fol}{tab}{der}}
      {\olfileid{pl}{tab}{der}}

\olsection{\usetoken{P}{tableau}}

\begin{explain}
অনুমিতি কী এবং অনুমান-বিধিগুলি কী, তা বলা হয়েছে। এগুলি থেকেই
আরোহীভাবে !!^{tableau}s তৈরি হয়: প্রতিটি !!{tableau} হয় এক বা
একাধিক অনুমিতি-সম্বলিত একটি শাখা, নয়তো কোনো শাখায় অনুমান-বিধির
একটি প্রয়োগের ফলে আগের কোনো !!a{tableau} থেকে উৎপন্ন হয়।
\end{explain}

\begin{defn}[!!^{tableau}]
$\sFmla{S_1}{!A_1}$, \dots, $\sFmla{S_n}{!A_n}$ অনুমিতিগুলির
(যেখানে প্রতিটি $S_i$ হয় $\True$, নয়তো~$\False$) জন্য
!!^a{tableau} হলো !!{signed formula}s-এর এমন একটি সসীম বৃক্ষ যা
নিচের শর্তগুলি পূরণ করে:
\begin{enumerate}
\item বৃক্ষের সর্বোচ্চ $n$টি !!{signed formula}s হলো
  $\sFmla{S_i}{!A_i}$, একটি আর-একটির নিচে।
\item বৃক্ষের যে !!{signed formula} অনুমিতি নয়, সেটি তার ওপরের
  শাখার কোনো !!a{signed formula}-এ অনুমান-বিধির সঠিক প্রয়োগের ফল।
\end{enumerate}
কোনো !!a{tableau}-এর একটি শাখায় $\sFmla{\True}{!A}$ ও
$\sFmla{\False}{!A}$ উভয়ই থাকলে এবং কেবল তখনই শাখাটি \emph{বদ্ধ};
অন্যথায় সেটি \emph{খোলা}। যে !!^a{tableau}-এর প্রতিটি শাখা বদ্ধ,
সেটি তার অনুমিতি-সমষ্টির জন্য একটি \emph{বদ্ধ !!{tableau}}। কোনো
!!a{tableau} বদ্ধ না হলে, অর্থাৎ তাতে অন্তত একটি খোলা শাখা থাকলে,
সেটি \emph{খোলা}।
\end{defn}

\begin{ex}
  যেকোনো অনুমিতি-সমষ্টি নিজেই !!a{tableau}, তবে সাধারণত সেটি বদ্ধ
  হবে না। (স্পষ্টতই, অনুমিতিগুলির মধ্যেই !!{signed formula}s-এর
  $\sFmla{\True}{!A}$ ও~$\sFmla{\False}{!A}$ রূপের কোনো জোড়া
  থাকলেই কেবল সেটি বদ্ধ।)

  কোনো !!a{tableau} (খোলা বা বদ্ধ) থেকে তার কোনো
  !!a{signed formula}~$!A$-তে একটি অনুমান-বিধি প্রয়োগ করে নতুন ও
  বড় ট্যাবলো পাওয়া যায়। বিধিটি যে~$!A$-এর আবির্ভাবে প্রয়োগ
  করা হয়েছে, সেই আবির্ভাব-সম্বলিত প্রতিটি শাখার শেষে এক বা একাধিক
  !!{signed formula}s যোগ করবে।

  যেমন, অনুমিতি $\sFmla{\True}{!A \land \lnot !A}$ বিবেচনা করো।
  শুধু সেই অনুমিতি নিয়ে গঠিত (খোলা) !!{tableau}-টি হলো:
  \begin{center}
    \begin{tableau}{}
      [\sFmla{\True}{\formula{A} \land \lnot \formula{A}}, just=\TAss]
    \end{tableau}{}
  \end{center}
  অনুমিতিটিতে $\TRule{\True}{\land}$ বিধি প্রয়োগ করে এটি থেকে নতুন
  !!{tableau} পাই। বিধিটি !!{tableau}-তে দুটি নতুন পংক্তি—
  $\sFmla{\True}{!A}$ ও $\sFmla{\True}{\lnot !A}$—যোগ করতে দেয়:
  \begin{center}
    \begin{tableau}{}
      [\sFmla{\True}{\formula{A} \land \lnot \formula{A}}, just=\TAss
        [\sFmla{\True}{\formula{A}}, just={\TRule{\True}{\land}[1]},
          [\sFmla{\True}{\lnot \formula{A}}, just={\TRule{\True}{\land}[1]}
          ]
        ]
      ]
    \end{tableau}{}
  \end{center}
  !!{tableau}s লেখার সময় ডানদিকে প্রয়োগ করা বিধিটি নথিবদ্ধ করি
  (যেমন, $\TRule{\True}{\land} 1$ বোঝায় যে সেই পংক্তির
  !!{signed formula}-টি পংক্তি~$1$-এর !!{signed formula}-এ
  $\TRule{\True}{\land}$ বিধি প্রয়োগের ফল)। নতুন !!{tableau}-টিতে
  এখন আরও !!{signed formula}s আছে, কিন্তু মাত্র একটিতে
  ($\sFmla{\True}{\lnot !A}$) বিধি প্রয়োগ করা যায়—এখানে
  $\TRule{\True}{\lnot}$ বিধি। ফলে নিচের বদ্ধ !!{tableau} পাই:
  \begin{center}
    \begin{tableau}{}
      [\sFmla{\True}{\formula{A} \land \lnot \formula{A}}, just=\TAss
        [\sFmla{\True}{\formula{A}}, just={\TRule{\True}{\land}[1]}
          [\sFmla{\True}{\lnot \formula{A}}, just={\TRule{\True}{\land}[1]}
            [\sFmla{\False}{\formula{A}}, just={\TRule{\True}{\lnot}[3]}, close]
          ]
        ]
      ]
    \end{tableau}{}
  \end{center}
\end{ex}

\end{document}
